\chapter{Introduction}
\label{chap:introduction}

\section{Parkinson's disease as a network disorder}

Parkinson's disease (PD) is a progressive neurodegenerative disorder characterized by both motor symptoms---including bradykinesia, rigidity, tremor, and postural instability---and non-motor symptoms involving cognition, mood, sleep, and autonomic function. Although degeneration of dopaminergic neurons in the substantia nigra is a defining pathological feature, the diversity of these manifestations cannot be explained by a single brain region. PD affects interactions among distributed cortical, subcortical, motor, cognitive, and limbic systems and is therefore increasingly understood as a network disorder \cite{Li-Chuan2019Graphtheory,Akbari2023Graphtheory}. There is still no well-established non-invasive biomarker for early diagnosis or progression monitoring, which limits the development and evaluation of disease-modifying interventions \cite{PPMI2011}.

Recent studies combining resting-state functional magnetic resonance imaging (rs-fMRI) with network neuroscience have provided valuable insights into altered brain organization in brain-related diseases. Rs-fMRI measures spontaneous fluctuations in the blood-oxygen-level-dependent signal while a participant is not performing an explicit task particularly make it more readily available and easier to share than t-fMRI as t-fMRI data varies due to the diverse task paradigms used in different studies. 

Network neuroscience complements rs-fMRI by providing tools to describe how these functional connections are organized across the whole brain. Using graph theory, brain regions are represented as nodes and functional connections as edges \cite{Bullmore2009Complex,Wang2010Graph}. Local graph measures, including degree, clustering coefficient, betweenness centrality, and eigenvector centrality, characterize the connectivity, segregation, and influence of individual regions. Global measures, such as average shortest path length and small-worldness, characterize large-scale integration and the balance between specialized and distributed processing. In this way, network neuroscience converts a high-dimensional rs-fMRI connectivity matrix into interpretable measures of regional and whole-network organization that can reveal distributed abnormalities missed by analyses of isolated regions.

Together, rs-fMRI and network neuroscience make it possible to investigate both where functional alterations occur and how they affect communication across the brain. Studies using this combination have reported PD-related alterations in cortico-striatal, sensorimotor, default-mode, and other intrinsic networks, as well as differences in degree, clustering, path length, efficiency, and centrality \cite{Jinping2017Graphtheory,Rutvi2021Graphtheory,Akbari2023Graphtheory}. However, reported findings are not always consistent, and graph estimates depend on preprocessing, connectivity estimation, parcellation, thresholding, software, and computational conditions. 

\section{Motivation for reproducibility analysis}
Reproducibility is essential to the accumulation of reliable scientific knowledge and is particularly important in neuroimaging, where results may ultimately inform clinical research. Open data, shared code, and detailed documentation improve transparency, but they do not guarantee that an analysis will produce consistent findings. Neuroimaging workflows contain many interconnected processing and statistical stages, so variation introduced by the data, analytical decisions, software tools, or computational environment can propagate to the final result.

One major challenge is \emph{analytical variability}, which arises from the many defensible choices available when constructing an fMRI workflow. These choices include participant selection, acquisition and preprocessing procedures, software package and version, parameter settings, statistical models, and correction for multiple comparisons \cite{carp2012plurality,li2024moving,germani2025hcp}. The Neuroimaging Analysis Replication and Prediction Study (NARPS) demonstrated the consequences of this flexibility: 70 teams analyzed the same fMRI dataset, yet no two teams used identical pipelines and, on average, 20\% of teams reported different conclusions \cite{botvinik2020variability}. Comparisons across neuroimaging tools have similarly shown that software packages, versions, and operating systems can alter statistical maps, tissue measurements, functional components, and graph properties \cite{Bhagwat2021,Glatard2015}. 

Graph-based analyses introduce additional decisions when defining nodes and edges: parcellation and spatial resolution determine the nodes, while Pearson correlation, partial correlation, mutual information, and other connectivity estimators can produce different network architectures \cite{Wang2010Graph,Fornito2010Scaling}. Reproducibility also depends on preprocessing choices such as global signal regression and smoothing, as well as graph sparsity or density; previous reviews have therefore reported reliability ranging from poor to moderate for some rs-fMRI graph measures \cite{Telesford2013Reproducibility}.

A second, less visible challenge is \emph{numerical variability}. Small differences in floating-point calculations can be introduced by hardware architectures, operating systems, compilers, mathematical libraries, or parallel execution. Although these perturbations begin at very small numerical scales, they can accumulate through long or nonlinear pipelines and produce meaningful differences in derived outcomes \cite{Glatard2015,Salari2021,Kiar2020,chatelain2024numerical}. Numerical variability may therefore contribute to the broader analytical variability observed across tools and infrastructures \cite{kiar2024experimental}. 

\section{Research problem}

Network neuroscience provides a powerful framework for studying the mechanisms underlying brain-related diseases. As analyses become increasingly computational, ensuring their numerical reliability has become a critical challenge.

While the numerical variability of structural imaging workflows has been investigated, with findings ranging from negligible to substantial, functional MRI (fMRI) pipelines and their derived graph measures remain underexplored. Without rigorous stability assessment, conclusions drawn from these measures may remain uncertain. 

This research addresses the gap by evaluating the numerical variability of graph measures of fMRI derived from the widely used fMRIPrep pipeline. It investigates how numerical perturbations influence the reliability of fMRI derived graph metrics, 
%across datasets and preprocessing configurations. 
providing insights about the robustness of network-based biomarkers and contributes to best practices for computational reproducibility in neuroimaging research.

The first major contribution of this thesis is systematic evaluation of the numerical variability of graph measures of functional connectivity derived from the widely used fMRIPrep pipeline and compared it to population variability. The resulting Numerical-Population Variability Ratio (\npvratio{}) values typically ranged from 0.02 to 0.19 for most graph metrics, indicating a measurable influence of numerical variability on network-derived outcomes. \npvratio{} values varied across brain regions, thresholding choices, and confound regression strategies. For results close to the significance threshold, propagation of the observed \npvratio{} values indicated an average 17.7\% probability that numerical variability would change a result from significant to non-significant, or vice versa.\ These findings highlight numerical variability as an important factor in functional network studies, particularly when examining subtle effects or working with small sample sizes.


This thesis proposal is organized around its main contributions, with each chapter addressing a distinct aspect of numerical uncertainty in numerical variability of graph measures of fMRI neuroimaging analyses. Before describing these contributions in detail, a theoretical background
on numerical variability is provided in Chapter~\ref{chap:background}. Chapter~\ref{chap:phase-one} describes the methods used to evaluate numerical variability in fMRI graph measures, including the datasets, preprocessing pipelines, and graph metrics analyzed. Chapter~\ref{chap:future-work} presents the proposed future work and subsequent research phases.
